% Solutions to the exercises of Chapter 5.
\begin{solutions}
  \item Let $F$ = faulty and $T$ = flagged. By the law of total probability $P(T) = 0.98\cdot0.01 + 0.03\cdot0.99 = 0.0098 + 0.0297 = 0.0395$, and by
        Bayes' rule
        \[
          P(F\mid T) = \frac{P(T\mid F)P(F)}{P(T)} = \frac{0.0098}{0.0395} = 0.248 .
        \]
        Only one flagged unit in four is actually faulty. The test is good, but faults are rare: the 3\% false-alarm rate applied to the 99\% of
        healthy units produces three times as many flags as the faulty units do. The prior matters as much as the likelihood.

  \item With $c = \cos30^\circ$ and $s = \sin30^\circ$, \cref{prop:affine} gives
        \[
          \Cov(y) = A\diag(1, 4)A\T = \begin{bmatrix} c^2 + 4s^2 & -3cs \\ -3cs & s^2 + 4c^2 \end{bmatrix}
          = \begin{bmatrix} 1.75 & -1.299 \\ -1.299 & 3.25 \end{bmatrix}.
        \]
        A rotation does not change the eigenvalues, still 1 and 4. The $1\sigma$ ellipse of $x$ has semi-axes 1 along $x_1$ and 2 along $x_2$; that
        of $y$ is the same ellipse rotated by $30^\circ$, its long axis (length 2) along $Ae_2 = (-\tfrac12, \tfrac{\sqrt3}{2})$, at $120^\circ$ from the
        first axis. The negative covariance says that $y_1$ and $y_2$ tend to have opposite signs, as the tilt of the ellipse shows.

  \item The difference is
        \[
          \Sigma_{aa} - \Sigma_{a|b} = \Sigma_{ab}\Sigma_{bb}\inv\Sigma_{ba} = \big(\Sigma_{bb}^{-1/2}\Sigma_{ba}\big)\T\big(\Sigma_{bb}^{-1/2}\Sigma_{ba}\big) \succeq 0,
        \]
        because $x\T\Sigma_{ab}\Sigma_{bb}\inv\Sigma_{ba}x = \|\Sigma_{bb}^{-1/2}\Sigma_{ba}x\|^2 \ge 0$ for every $x$ (with $\Sigma_{bb}^{-1/2}$ any square root of
        $\Sigma_{bb}\inv \succ 0$). Equality holds for all $x$ exactly when $\Sigma_{ba}x = 0$ for all $x$, i.e.\ when $\Sigma_{ab} = 0$: $a$ and $b$ are uncorrelated,
        hence (being jointly Gaussian) independent, and observing $b$ says nothing about $a$.

  \item Order the variables as $(b, a)$, so that $\Lambda$ becomes $\begin{bmatrix}\Lambda_{bb} & \Lambda_{ba}\\ \Lambda_{ab} & \Lambda_{aa}\end{bmatrix}$ and its inverse is
        $\begin{bmatrix}\Sigma_{bb} & \Sigma_{ba}\\ \Sigma_{ab} & \Sigma_{aa}\end{bmatrix}$. By \cref{prop:block-inverse}, with the top-left block $\Lambda_{bb}$ playing the
        role of $A$, the bottom-right block of the inverse is the inverse of the Schur complement:
        \[
          \Sigma_{aa} = \big(\Lambda_{aa} - \Lambda_{ab}\Lambda_{bb}\inv\Lambda_{ba}\big)\inv
          \quad\Longleftrightarrow\quad \Lambda_{aa} - \Lambda_{ab}\Lambda_{bb}\inv\Lambda_{ba} = \Sigma_{aa}\inv .
        \]
        The marginal of $a$ is $\N(\mu_a, \Sigma_{aa})$, whose information matrix is $\Sigma_{aa}\inv$: in information form, marginalization is this Schur
        complement. (The same proposition gives the marginal information vector $\eta_a - \Lambda_{ab}\Lambda_{bb}\inv\eta_b$.)

  \item Up to terms that do not involve $x$, the exponent is
        \[
          -\frac{(x - z_1)^2}{2\sigma_1^2} - \frac{(x - z_2)^2}{2\sigma_2^2}
          = -\frac12\Big(\frac{1}{\sigma_1^2} + \frac{1}{\sigma_2^2}\Big)x^2 + \Big(\frac{z_1}{\sigma_1^2} + \frac{z_2}{\sigma_2^2}\Big)x + \text{const}
          = -\frac{(x - \hat x)^2}{2\sigma^2} + \text{const},
        \]
        where matching the coefficient of $x^2$ gives $1/\sigma^2 = 1/\sigma_1^2 + 1/\sigma_2^2$ and matching that of $x$ gives
        $\hat x/\sigma^2 = z_1/\sigma_1^2 + z_2/\sigma_2^2$. A density proportional to $\exp(-(x - \hat x)^2/2\sigma^2)$ is $\N(\hat x, \sigma^2)$. Finally, with
        $\sigma^2 = \sigma_1^2\sigma_2^2/(\sigma_1^2 + \sigma_2^2)$,
        $\hat x = (\sigma_2^2z_1 + \sigma_1^2z_2)/(\sigma_1^2 + \sigma_2^2) = z_1 + \sigma_1^2(z_2 - z_1)/(\sigma_1^2 + \sigma_2^2)$, the second form of \eqref{eq:fusion}.

  \item The cost $J(x) = (x - \mu)\T P\inv(x - \mu) + (z - Hx)\T R\inv(z - Hx)$ has gradient $2P\inv(x - \mu) - 2H\T R\inv(z - Hx)$. Setting it to zero,
        \[
          \big(P\inv + H\T R\inv H\big)\hat x = P\inv\mu + H\T R\inv z .
        \]
        Write $\hat x = \mu + \delta$ and subtract $(P\inv + H\T R\inv H)\mu$ from both sides: $(P\inv + H\T R\inv H)\,\delta = H\T R\inv(z - H\mu)$, so
        $\hat x = \mu + K(z - H\mu)$ with $K = (P\inv + H\T R\inv H)\inv H\T R\inv$. The Hessian $2(P\inv + H\T R\inv H) \succ 0$ confirms a minimum. By the push-through
        identity $(P\inv + H\T R\inv H)\inv H\T R\inv = PH\T(HPH\T + R)\inv$ this is the Kalman gain of \cref{ch:kf}.

  \item Exactly, $\E[y] = \E[x^2] = \sigma^2$ (and $\Var(y) = \E[x^4] - \sigma^4 = 2\sigma^4$). Linearizing $f(x) = x^2$ at the mean $\mu = 0$ gives $f(0) = 0$ and $J = 2\mu = 0$:
        the predicted mean is $0$ and the predicted variance is $0$, a deterministic $y$. The linearization sees a function that is flat at the
        expansion point and misses the curvature entirely. The second-order correction $\tfrac12f''(\mu)\sigma^2 = \sigma^2$ of \cref{sec:linearization} accounts
        for the whole discrepancy (exactly, because $f$ is quadratic). This is the worst case for an EKF: linearizing at a stationary point of
        the measurement function.

  \item The squared Mahalanobis distance of a 2D Gaussian is $\chi^2_2$, whose density $\tfrac12e^{-\gamma/2}$ integrates to $P(\chi^2_2 \le c) = 1 - e^{-c/2}$.
        The $k\sigma$ ellipse is $d_M^2 = k^2$:
        \[
          P(1\sigma) = 1 - e^{-1/2} = 0.393, \qquad P(2\sigma) = 1 - e^{-2} = 0.865, \qquad P(3\sigma) = 1 - e^{-9/2} = 0.989 .
        \]
        Compare $0.683$, $0.954$ and $0.997$ in one dimension (and $0.199$, $0.739$, $0.971$ in three): the more dimensions, the more likely a sample lies outside
        the ``one-sigma'' region, because it must be close to the mean in every direction at once.

  \item With $\varphi(s) = e^{-s^2/2}/\sqrt{2\pi}$ we have $\varphi'(s) = -s\varphi(s)$. Integrating by parts with $u = s^3$ and $\dd v = s\varphi\,\dd s = -\dd\varphi$,
        \[
          \E[s^4] = \int_{-\infty}^{\infty}s^3\cdot s\varphi(s)\,\dd s = \big[-s^3\varphi(s)\big]_{-\infty}^{\infty} + 3\int_{-\infty}^{\infty}s^2\varphi(s)\,\dd s = 0 + 3\,\E[s^2] .
        \]
        The same step with $u = s$ gives $\E[s^2] = [-s\varphi]_{-\infty}^{\infty} + \int\varphi = 1$. Hence $\E[s^4] = 3$ and $\Var(s^2) = 3 - 1 = 2$, the variance used for $\chi^2_k$.

  \item Under the null hypothesis that the predicted variance is right, $99\,s^2/\sigma^2 \sim \chi^2_{99}$, so the ratio $s^2/\sigma^2$ lies in
        $[\chi^2_{0.025}(99), \chi^2_{0.975}(99)]/99 = [0.741, 1.297]$ with probability 95\% (\cref{app:tables}). The observed ratio $1.25$ is inside the
        interval: \emph{not} significant at the 5\% level (the two-sided $p$-value is about $0.09$). With 100 runs, a variance can be off by 25\%
        without the test noticing; the expected ratio $1.25$ would leave the interval only from about 140 runs, and a 10\% error needs
        roughly a thousand. Monte Carlo variance checks need many runs (Exercise 17.6).
\end{solutions}
